\documentclass[11pt]{article}
\usepackage[margin=1in]{geometry}
\usepackage{amsmath,amssymb,amsthm}
\usepackage{booktabs}
\usepackage{array}
\usepackage[T1]{fontenc}
\usepackage{lmodern}
\usepackage{microtype}

\newcommand{\Aperp}{A_{\perp}}
\newcommand{\Gperp}{G_{\perp}}
\newcommand{\GF}{G_F}
\newcommand{\Cperp}{C_{\perp}}
\newcommand{\Wperp}{W_{\perp}}
\newcommand{\mF}{m_F}
\newtheorem{theorem}{Theorem}
\newtheorem{lemma}{Lemma}

\title{A Projected Frobenius Envelope Gap for the Literal V59 Output Lane}
\author{Liang Wang\\
School of Mathematics and Statistics\\
Huazhong University of Science and Technology (HUST), Wuhan, China}
\date{27 August 2026}

\begin{document}
\maketitle

\begin{abstract}
TPC-273 showed that a finite correlation margin can move across quantitative
bands when a declared cutoff is changed.  We next test the standard
cancellation-free response: apply a Frobenius norm envelope after the same
three-block Haar projection.  On the locked literal V59 operator, let
\(\Aperp=(I-P_3)A\), \(g= A\beta\), and
\(\Gperp=\|\Aperp\beta\|_2^2\).  We prove the exact finite inequality
\(\Gperp\leq\|\Aperp\|_F^2\|\beta\|_2^2=:\GF\).  Exact rational construction and
independent replay on six registered growing-cutoff scales and two kernel
exponents certify \(\GF/\Gperp>50\) and the envelope margin proxy
\(\mF^2<1/64\) on all 12 rows.  This is a quantitative, scoped insufficiency
of a norm-only output proof; it is not an upper bound on the actual margin,
an asymptotic counterexample, or a twin-prime theorem.
\end{abstract}

\section{Question and claim ceiling}

The residual route separates a signed scalar \(\Cperp\) from two norm lanes.  If
\(\Wperp=\|(I-P_3)w\|_2^2\) and \(\Gperp=\|(I-P_3)g\|_2^2\), then the correlation
margin is
\[
 m^2=\frac{|\Cperp|^2}{\Wperp\Gperp}.
\]
TPC-272 converted a signed saving and a margin loss into an endpoint budget;
TPC-273 then demonstrated finite cutoff sensitivity.  The present paper asks
whether a simple output estimate can pay that missing margin cost.

Our claim ceiling is
\[
\texttt{PROVED\_EXACT\_FINITE} \quad+\quad
\texttt{NUMERICALLY\_CERTIFIED\_FINITE},
\]
with a scoped `\texttt{INSUFFICIENT\_SCOPED}` verdict for the tested
cancellation-free route.  No fixed-power credit is assigned.

\section{The locked finite operator}

For a physical scale \(N\), put \(L=N/2\), and let \(\beta\) be the exact
source vector supplied by the released TPC-268 engine.  On the prime shell
\(Q<q\leq 2Q\), its literal matrix is
\[
 A_{u,t}=\mathbf 1_{u\ne t}\sum_q qK_{H,s}(u-t)
 \mathbf 1_{q\nmid u}\mathbf 1_{q\nmid t}
 \left(\mathbf 1_{u\equiv t\pmod q}-\frac1{q-1}\right),
\]
where \(K_{H,s}(h)=(1+(h/H)^2)^{-s}\).  The diagonal deletion, unit masks,
prime shell, beta, and four consecutive blocks are frozen; only the registered
kernel exponent \(s\in\{1,2\}\) is paired as a control.

The six growing-cutoff rows are
\[
 (N,H,Q)=(64,15,4),(96,20,5),(128,24,5),
 (192,32,6),(256,38,6),(384,50,7),
\]
with comparison cutoff \(z_N=4,4,4,5,5,5\), respectively.  Let \(P_3\)
denote the orthogonal projection onto the three declared four-block Haar
contrasts and set
\[
 \Aperp=(I-P_3)A,\qquad g_{\perp}=\Aperp\beta,qquad
 \Gperp=\|g_{\perp}\|_2^2.
\]
All entries of \(A\), \(P_3\), and \(\beta\) are rational on these rows.

\section{The projected Frobenius theorem}

\begin{theorem}[projected Frobenius envelope]
For every finite real or complex matrix \(B\) and vector \(v\),
\[
 \|Bv\|_2^2\leq \|B\|_F^2\|v\|_2^2.
\]
Consequently, on every registered row,
\[
 \Gperp\leq \GF:=\|\Aperp\|_F^2\|\beta\|_2^2.
\]
\end{theorem}

\begin{proof}
For each row \(i\), Cauchy--Schwarz gives
\[
 \left|\sum_jB_{ij}v_j\right|^2
 \leq\left(\sum_j|B_{ij}|^2\right)\left(\sum_j|v_j|^2\right).
\]
Summing over \(i\) yields the first inequality.  Substitution of
\(B=\Aperp\) and \(v=\beta\) gives the second.  In particular, the theorem
does not use an independence or probabilistic assumption; its cost is that it
discards cancellation between the columns of \(\Aperp\).
\end{proof}

For positive finite lanes define the conservative proxy
\[
 \mF^2=\frac{|\Cperp|^2}{\Wperp\GF}.
\]
The theorem implies \(\mF^2\leq m^2\).  Thus a small \(\mF\) is not a
counterexample to a large actual margin: it only says that this envelope cannot
certify one.

\section{Exact finite audit}

The producer constructs every matrix entry and projection coefficient as a
`\texttt{Fraction}`.  It checks \(A\beta=g\) against the parent engine, computes
\(\|\Aperp\|_F^2\), \(\|\beta\|_2^2\), and \(\GF\) exactly, and transfers only
the parent outward intervals for \(\Cperp\), \(\Wperp\), and \(\Gperp\).  A
separate checker repeats the matrix construction without importing the
producer.  A five-mutation stress audit rejects altered theorem counts,
envelope values, threshold metadata, operator status, and fixed-power credit.

\begin{table}[t]
\centering
\scriptsize
\begin{tabular}{c c c r r c}
\toprule
 $N$ & $s$ & $z_N$ & lower gap $\GF/\Gperp$ & upper $\mF^2$ & phase\\
\midrule
 64  & 1 & 4 & 63.25 & 0.001183 & negative\\
 64  & 2 & 4 & 51.35 & 0.001197 & negative\\
 96  & 1 & 4 & 128.83 & 0.000698 & negative\\
 128 & 1 & 4 & 237.09 & 0.000066 & negative\\
 192 & 1 & 5 & 174.15 & 0.000001 & negative\\
 256 & 2 & 5 & 247.09 & 0.000001 & positive\\
 384 & 1 & 5 & 479.52 & 0.000003 & negative\\
\bottomrule
\end{tabular}
\caption{Selected outward finite bounds.  The complete certificate has 12
rows; displayed decimals are only for readability.}
\end{table}

\begin{theorem}[registered finite gap]
On all 12 rows, the exact envelope satisfies
\(\GF/\Gperp>50\) and the transferred conservative proxy satisfies
\(\mF^2<1/64\).  The phase intervals are separated from zero: 11 are on the
negative real axis and one is on the positive real axis.
\end{theorem}

\begin{proof}
The six scale rows are paired with \(s=1,2\), giving 12 distinct keys.  The
certificate stores the exact rational Frobenius energy and beta energy, then
divides the resulting exact \(\GF\) by the parent outward interval for
\(\Gperp\).  It also divides the squared scalar interval by the positive
\(\Wperp\) interval and exact \(\GF\).  The lower and upper endpoints pass the
strict thresholds on every row.  Independent reconstruction reproduces all
exact fields and interval transfers.  The phase count is a direct sign count
of the stored scalar intervals.
\end{proof}

The gap is not a claim about a limiting sequence.  It is a finite diagnostic
that the rowwise norm envelope can be much looser than the actual output for
the locked source.  The exponent control is similarly descriptive: the six
envelope ratios for the transition \(s=1\to2\) range from about \(0.62\) to
\(0.63\), while the actual output ratios remain distinct, so the envelope does
not reproduce the source cancellation.

\section{Route evaluation and limits}

The exact inequality is a reusable structural lemma, and the finite gap is a
new numerical certificate.  Its negative conclusion is deliberately narrow:
the cancellation-free projected Frobenius route is insufficient on the
registered interface.  It does not refute a signed estimate using the actual
coefficient phases.  In particular, the following gates remain open:
\[
 \begin{gathered}
 \texttt{SOURCE\_LEVEL\_OUTPUT\_BOUND}=\texttt{OPEN\_ASYMPTOTIC}\\
 \texttt{SIGNED\_FOUR\_PACKET\_REASSEMBLY}=\texttt{OPEN}\\
 \texttt{ARITHMETIC\_L2}=\texttt{NONE}.
 \end{gathered}
\]
The finite rows also contain one positive phase, which is retained rather than
removed by a sign convention.  No row count is converted into a power saving,
and the required strict \(1/400\) endpoint payment remains unpaid.

\section{Conclusion}

TPC-274 closes one low-cost proof attempt: after the physical projection, a
valid Frobenius envelope loses more than a factor of 50 on every registered
row and yields a margin proxy below \(1/8\) everywhere.  The natural next
question is signed output reassembly, where the coefficient phases are kept
instead of discarded.  That question is not answered here; the present result
is a finite, reproducible obstruction to the norm-only shortcut.

\begin{thebibliography}{9}
\bibitem{HW}
G. H. Hardy and E. M. Wright, \emph{An Introduction to the Theory of Numbers},
6th ed., Oxford University Press, 2008.
\bibitem{TPC268}
L. Wang, ``Finite Cutoff-Sensitivity Obstruction for a V59 Residual,'' TPC-268
project release, 2026.
\bibitem{TPC273}
L. Wang, ``A Finite Margin-Stability Matrix for the Literal V59 Residual,''
TPC-273 project release, 2026.
\end{thebibliography}

\end{document}
